The assembly itemises its joints in its statement; each joint carries its own provenance there.

\emph{Premise (I) (\claimref{hawking.premise-i-interface-identification}), the interface identification:} the horizon-sector state feeds the exterior modes in the standard way, i.e.\ the occupancy of each near-horizon channel is carried to the corresponding asymptotic channel by the exterior propagation. This is the same assumption every derivation of Hawking flux makes, from \cite{Hawking1975} onward, and it is this assembly's one non-theorem joint. Under (I), the emitted flux in each channel $(\omega, \ell, m, s)$ equals (with $\mathcal{N}$ the count of emitted quanta)
\[
\frac{d\mathcal{N}}{dt\,d\omega} \;=\; \frac{1}{2\pi}\,\Gamma_{s\ell m}(\omega)\; n_\pm(\omega), \qquad
n_\pm(\omega) = \frac{1}{e^{\omega/\TH} \pm 1}
\]
(the upper sign for the fermionic species, the lower for the bosonic sector), with $\Gamma_{s\ell m}$ the classical transmission coefficient through the Regge-Wheeler or Teukolsky barrier. The joints: (i) the barriers and their transmission coefficients are properties of the exterior metric, and by Fact \claimref{hawking.imported-exterior-exactly-general-relativity} the framework's exterior description \emph{is} classical general relativity on the invertible region, exactly and not asymptotically, so $\Gamma_{s\ell m}$ is inherited verbatim, mode by mode \cite{ReggeWheeler1957, Teukolsky1973}; (ii) the occupancies at the horizon are Theorem \claimref{hawking.i-state-continuum} for the fermionic species and Fact \claimref{hawking.imported-s-wave-bosonic-sector}'s (Q)-conditional half for the bosonic sector, both at exactly $\TH$ with the temperature's coefficient fixed geometrically (Fact \claimref{hawking.imported-modular-descent-premises-geometric-normalisation}); (iii) given (I), the displayed formula is the definition of filtered thermal flux: the framework's flux prediction equals Hawking's \cite{Hawking1975, Page1976}, greybody structure included.
